\chapter*{Úvod} % chapter* je necislovana kapitola
\addcontentsline{toc}{chapter}{Úvod} % rucne pridanie do obsahu
\markboth{Úvod}{Úvod} % vyriesenie hlaviciek

Predmet Logika pre informatikou sa zaoberá logikou prvého rádu. Pri výučbe predmetu študenti používajú webový interaktívny logický pracovný zošit Logic Workbook, ktorý vznikol ako súčasť bakalárskej práce Mateja Moka \cite{logicworkbook}. Jednou z tém, ktoré sa vyučujú počas predmetu je dôkaz tvrdení pomocou rezolvencie, pri čom sa využíva editor rezolvenčných dôkazov, ktorý bol vytvorený Norbertom Juríkom \cite{rezolvovac}. Na riešenie rezolvencie je potrebný vstup (klauzálna teória), ktorý je vytvorený postupným aplikovaním ekvivalentných úprav a skolemizácie na počiatočné formuly. Logic Workbook neobsahuje nástroj, ktorý by kontroloval jednotlivé úpravy. Z toho vyplýva, že študenti musia robiť tento postup ručne, pričom ani študenti, ani vyučujúci nemajú ľahký spôsob kontroly správnosti úprav.

Cieľom práce je vytvoriť novú aplikáciu, ktorá bude integrovaná do Logic Workbooku. Aplikácia bude slúžiť na kontrolu správneho prevodu formuly do konjunktívnej normálnej formy. Keďže je určená ako učebná pomôcka, aplikácia bude kontrolovať iba správnosť jednotlivých ekvivalentných úprav a skolemizácie. Ak by prevod bol príliš automatizovaný mohla by nastať situácia, kde študenti nebudú rozumieť postupu, ale iba dosiahnu správny výsledok spôsobom pokus-omyl. Preto samotný postup prevodu bude ponechaný na študentov. Aplikácia umožní vybrať jeden druh ekvivalentnej úpravy, ktorá sa v danom kroku uplatní buď na počiatočnú formulu alebo formulu vytvorenú predošlou ekvivalentnou úpravou. Aplikácia v každom kroku skontroluje či sú dve formuly ekvivalentné podľa vybranej úpravy. Pomocou postupností týchto úprav, študenti budú môcť skontrolovať svoj postup. Keďže študenti často potrebujú previesť viacero formúl do konjunktívnej normálnej formy, budú môcť aplikácii zadať viacero počiatočných formúl. Každá počiatočná formula bude začiatkom osobitnej postupnosti ekvivalentných úprav. Keďže aplikácia bude integrovaná do logického pracovného zošit, jej používateľské rozhranie bude s ním kompatibilné.

V prvej kapitole zhrnieme teoretické znalosti.
V druhej kapitole krátko opíšeme existujúci ekosystém a zhrnieme technológie, ktoré sme použili pri tvorbe aplikácie.
Potom v tretej kapitole sa budeme venovať návrhu aplikácie, jej internému stavu, základnému návrhu kontroly ekvivalentných úprav a tiež zhrnieme požiadavky pre aplikáciu.
V štvrtej kapitole bližšie popíšeme samotnú implementáciu aplikácie, jej súborovú štruktúru a používateľské rozhranie. Tiež opíšeme viacero špecifických implementácií kontroly správnosti úprav. 
Nakoniec v piatej a poslednej kapitole popíšeme ako sme aplikáciu testovali a zhrnieme výsledky testovania.